\section{Q-learning 调度算法设计}

在上一章中，本文将操作系统任务调度问题抽象为一个马尔可夫决策过程（Markov Decision Process, MDP）。在此基础上，本章进一步设计基于强化学习的调度算法。本文选择 Q-learning 作为核心学习算法，并结合操作系统调度场景，对状态空间、动作空间以及奖励函数进行具体设计，从而构建适用于 xv6 操作系统的强化学习调度模型。

\subsection{Q-learning 算法原理}

Q-learning 是一种经典的无模型强化学习算法，其目标是通过不断与环境交互学习状态—动作价值函数（Q-value），从而获得最优策略。

在强化学习框架中，智能体（Agent）在状态 $s$ 下选择动作 $a$，环境返回奖励 $r$，并转移到新状态 $s'$。Q-learning 通过不断更新 Q 值函数来逼近最优策略。

Q-learning 的核心更新公式如下：

\begin{equation}
Q(s,a) \leftarrow Q(s,a) + \alpha \left(r + \gamma \max_{a'} Q(s',a') - Q(s,a)\right)
\end{equation}

其中：

\begin{itemize}
\item $Q(s,a)$ 表示在状态 $s$ 下执行动作 $a$ 的价值估计
\item $\alpha$ 为学习率（learning rate），用于控制更新幅度
\item $\gamma$ 为折扣因子（discount factor），用于衡量未来奖励的重要程度
\item $r$ 为当前动作获得的即时奖励
\item $\max Q(s',a')$ 表示下一状态下可能获得的最大 Q 值
\end{itemize}

通过不断迭代更新 Q 表，算法能够逐渐学习到在不同状态下最优的动作选择。

为了在探索（exploration）与利用（exploitation）之间取得平衡，本文采用 $\varepsilon$-greedy 策略进行动作选择：

\begin{itemize}
\item 以概率 $\varepsilon$ 随机选择动作（探索）
\item 以概率 $1-\varepsilon$ 选择当前 Q 值最大的动作（利用）
\end{itemize}

该策略能够避免算法过早陷入局部最优，同时保证学习过程的稳定性。

在操作系统调度问题中，调度器可以视为强化学习中的智能体，而系统运行环境则构成强化学习环境。通过不断进行调度决策并更新 Q 表，调度器能够逐步学习到更加合理的任务调度策略。

\subsection{状态空间设计}

在强化学习调度模型中，状态空间用于描述当前系统的运行情况。合理的状态设计能够反映系统负载和进程运行特征，从而帮助算法做出更加有效的调度决策。

结合 xv6 操作系统的实现，本文选择以下指标作为状态特征：

\begin{itemize}
\item 进程等待时间（wait time）
\item 进程运行时间（run time）
\item 调度次数（schedule count）
\item 就绪队列长度（runnable processes）
\end{itemize}

上述信息均可以通过扩展 xv6 的 \texttt{struct proc} 结构体获得。

由于 Q-learning 算法通常适用于离散状态空间，因此需要对连续变量进行离散化处理。例如，可以将等待时间划分为若干区间：

\begin{itemize}
\item $0 \sim 10$ ticks
\item $10 \sim 50$ ticks
\item $>50$ ticks
\end{itemize}

类似地，可以对运行时间和调度次数进行分段，从而构建有限大小的状态空间。通过这种方式，可以将系统状态表示为：

\begin{equation}
s = (wait\_level, run\_level, queue\_length)
\end{equation}

该状态能够较好地反映进程在系统中的运行情况，为调度决策提供依据。

\subsection{动作空间设计}

动作空间表示在给定状态下调度器可以执行的操作。在操作系统调度问题中，调度器的主要任务是从就绪队列中选择一个进程运行。

因此本文将动作空间定义为：

\begin{equation}
A = \{选择某个 RUNNABLE 进程\}
\end{equation}

假设当前系统中存在 $N$ 个可运行进程，则动作空间大小为：

\begin{equation}
|A| = N
\end{equation}

在实际实现中，调度器会遍历 \texttt{proc[]} 表，找出所有状态为 RUNNABLE 的进程，并根据 Q 值选择其中最优的进程运行。

在 $\varepsilon$-greedy 策略下，调度器在每次调度时：

\begin{itemize}
\item 以概率 $\varepsilon$ 随机选择一个 RUNNABLE 进程
\item 以概率 $1-\varepsilon$ 选择 Q 值最大的进程运行
\end{itemize}

这种策略既保证了对未知状态的探索，又能够利用已有经验优化调度效果。

\subsection{奖励函数设计}

在基于强化学习的调度框架中，奖励函数用于刻画调度策略的优化目标，并指导智能体在不同状态下进行决策。由于操作系统调度目标（如平均等待时间、响应时间等）难以通过单步决策直接精确建模，本文采用启发式奖励函数进行近似刻画。

设当前状态为 $s$，执行动作 $a$ 后转移到下一状态 $s'$，奖励函数定义如下：

\begin{itemize}
    \item 当进程执行结束（终止状态）时，给予正奖励 $+5$，以鼓励调度策略提高任务完成效率；
    \item 若系统最大等待时间降低（即等待时间桶编号减小），给予正奖励 $+2$；
    \item 若系统最大等待时间增加，则给予负奖励 $-2$；
    \item 若状态无显著变化，则奖励为 $0$。
\end{itemize}

形式化表示为：

\[
R(s,a,s') =
\begin{cases}
+5, & \text{若达到终止状态} \\
+2, & \text{若等待时间下降} \\
-2, & \text{若等待时间上升} \\
0,  & \text{其他情况}
\end{cases}
\]

该奖励函数从两个方面刻画调度性能：一是通过终止奖励鼓励系统尽快完成任务；二是通过等待时间变化反映调度过程中的公平性与响应性。该设计属于典型的奖励塑形（reward shaping）方法，在保证实现复杂度较低的同时，为强化学习算法提供了有效的优化方向。

需要指出的是，该奖励函数为对真实调度目标的启发式近似，其主要作用在于引导策略学习，而非严格刻画全局最优性能指标。
\subsection{本章小结}

本章介绍了基于 Q-learning 的任务调度算法设计方法。首先介绍了 Q-learning 算法的基本原理及其更新公式，并采用 $\varepsilon$-greedy 策略实现探索与利用之间的平衡。随后结合操作系统调度场景，对状态空间、动作空间以及奖励函数进行了具体设计。在此基础上构建了适用于 xv6 操作系统的强化学习调度模型，为下一章的系统实现提供理论基础。